\documentclass[a4paper,12pt]{article}

% -------------------------------------------------
% Кодировка и язык
% -------------------------------------------------
\usepackage[T2A]{fontenc}
\usepackage[utf8]{inputenc}
\usepackage[russian]{babel}

% -------------------------------------------------
% Математика и типографика
% -------------------------------------------------
\usepackage{amsmath,amssymb,mathtools}
\usepackage{helvet}
\renewcommand{\familydefault}{\sfdefault}
\usepackage{icomma}
\usepackage[a4paper,margin=2.2cm]{geometry}
\usepackage{microtype}
\usepackage{setspace}
\onehalfspacing

% -------------------------------------------------
% Графика и таблицы
% -------------------------------------------------
\usepackage{tikz}
\usetikzlibrary{calc}
\usepackage{tabularx,booktabs}
\usepackage{array}

% -------------------------------------------------
% Списки и служебное оформление
% -------------------------------------------------
\usepackage{enumitem}
\usepackage{parskip}
\usepackage{fancyhdr}
\usepackage{xcolor}

% -------------------------------------------------
% Палитра
% -------------------------------------------------
\definecolor{accent}{HTML}{008080}
\definecolor{darkaccent}{HTML}{004D4D}
\definecolor{textgray}{HTML}{333333}
\definecolor{softgray}{HTML}{E9EEEE}

\color{textgray}

% -------------------------------------------------
% Параметры списков
% -------------------------------------------------
\setlist[itemize]{
  leftmargin=*,
  itemsep=0.35em,
  topsep=0.45em
}

\setlist[enumerate]{
  leftmargin=*,
  itemsep=0.45em,
  topsep=0.45em
}

% -------------------------------------------------
% Колонтитулы
% -------------------------------------------------
\pagestyle{fancy}
\fancyhf{}
\setlength{\headheight}{15pt}
\lhead{\small\color{darkaccent}Математика · 6 класс}
\rhead{\small Ксения Васильченко}
\cfoot{\small\thepage}
\renewcommand{\headrulewidth}{0.4pt}
\renewcommand{\headrule}{
  \hbox to\headwidth{\color{accent!55}\leaders\hrule height \headrulewidth\hfill}
}

\raggedbottom

% -------------------------------------------------
% Пользовательские команды
% -------------------------------------------------
\newcommand{\sectiontitle}[1]{%
  \section*{\color{darkaccent}#1}%
}

\newcommand{\important}[1]{%
  \textbf{\color{darkaccent}#1}%
}

\newcommand{\checkmarkitem}{%
  \item[\(\square\)]%
}

\begin{document}

% =================================================
% ТИТУЛЬНЫЙ ЛИСТ
% =================================================

\begin{titlepage}
\thispagestyle{empty}

\begin{center}

\vspace*{2.4cm}

{\Large\bfseries\color{accent} ЧЕК-ЛИСТ}

\vspace{1.1cm}

{\Huge\bfseries
Десятичные дроби
}

\vspace{0.45cm}

{\Large
сравнение и арифметические действия
}

\vspace{0.25cm}

{\Large
среднее арифметическое и простейшие уравнения
}

\vspace{1.7cm}

{\large
Математика · 6 класс
}

\vspace{2.0cm}

{\large
Лёвин Артём Александрович\\[0.4em]
\textbf{эксклюзивно для Ксении Васильченко}
}

\vspace{1.5cm}

{\large \today}

\end{center}

\vfill

\begin{center}
\begin{tikzpicture}[x=1cm,y=1cm]
  \draw[
    accent!55,
    line width=1.1pt
  ]
  (0,0)
  .. controls (2.5,0.55) and (4.5,0.55) .. (7,0)
  .. controls (9.5,-0.55) and (11.5,-0.55) .. (14,0);
\end{tikzpicture}
\end{center}

\vspace{0.6cm}

\end{titlepage}

% =================================================
% ОСНОВНОЙ ЧЕК-ЛИСТ
% =================================================

\sectiontitle{1. Письменные вычисления и порядок действий}

\important{При письменном сложении и вычитании натуральных чисел}
цифры располагают строго по разрядам:
единицы под единицами, десятки под десятками, сотни под сотнями.

Например,
\[
8000-3746=4254.
\]

Проверка вычитания выполняется обратным действием:
\[
4254+3746=8000.
\]

\important{Порядок действий:}

\begin{enumerate}
  \item действия в скобках;
  \item умножение и деление;
  \item сложение и вычитание.
\end{enumerate}

Например,
\[
720-6\cdot(20+15).
\]

Сначала:
\[
20+15=35,
\]
затем:
\[
6\cdot35=210,
\]
после этого:
\[
720-210=510.
\]

\important{Самопроверка.}
Перед окончательной записью ответа полезно ещё раз проверить,
какое действие должно выполняться первым.

% =================================================

\sectiontitle{2. Десятичные дроби: запись и чтение}

Десятичная дробь состоит из \important{целой} и
\important{дробной} частей, разделённых запятой.

\[
4{,}8
\]

читается: \textit{четыре целых восемь десятых}.

Количество цифр после запятой определяет разряд последней цифры:

\[
\begin{aligned}
4{,}8   &\quad \text{--- четыре целых восемь десятых},\\
4{,}08  &\quad \text{--- четыре целых восемь сотых},\\
4{,}008 &\quad \text{--- четыре целых восемь тысячных}.
\end{aligned}
\]

\subsection*{\color{accent}Нули справа от дробной части}

Нули, приписанные \important{справа в конце дробной части},
значение числа не изменяют:

\[
4{,}8=4{,}80=4{,}800.
\]

Поэтому:

\[
5{,}3000=5{,}3.
\]

Однако нули внутри записи числа удалять нельзя:

\[
4{,}08\ne4{,}8,
\]

\[
4{,}008\ne4{,}08.
\]

Например,
\[
4{,}08>4{,}008,
\]
поскольку после одинаковых целых частей и одинакового разряда
десятых получаем:
\[
8>0
\]
в разряде сотых.

% =================================================

\sectiontitle{3. Сравнение десятичных дробей}

Для сравнения двух десятичных дробей удобно действовать по алгоритму.

\begin{enumerate}
  \item Сравните целые части.
  \item Если целые части равны, при необходимости допишите справа нули.
  \item Сравнивайте цифры после запятой слева направо:
        десятые, затем сотые, тысячные и так далее.
  \item Первая различающаяся цифра определяет результат сравнения.
\end{enumerate}

Например:
\[
1{,}23=1{,}230,
\]
поэтому
\[
1{,}230>1{,}215.
\]

Действительно:

\[
\begin{array}{c@{\qquad}c}
1{,}230 & 1{,}215\\
\phantom{1,}2 & \phantom{1,}2\\
\phantom{1,2}3 & \phantom{1,2}1
\end{array}
\]

В разряде сотых:
\[
3>1.
\]

Если различаются целые части, достаточно сначала сравнить именно их:

\[
12{,}345>1{,}9999.
\]

\important{Ключевая идея:}
сравниваются цифры одного и того же разряда.

% =================================================

\sectiontitle{4. Сложение десятичных дробей}

\important{Главное правило: запятая должна находиться под запятой.}

Если количество цифр после запятой различается,
справа можно дописать нули.

Например:
\[
2{,}41+1{,}159.
\]

Запишем:
\[
2{,}410+1{,}159.
\]

Тогда:

\[
\begin{array}{r}
  2{,}410\\
+ 1{,}159\\
\hline
  3{,}569
\end{array}
\]

Следовательно,
\[
\boxed{2{,}41+1{,}159=3{,}569}.
\]

\important{Алгоритм:}

\begin{enumerate}
  \item расположите числа запятая под запятой;
  \item при необходимости допишите нули справа;
  \item складывайте справа налево;
  \item запятую в ответе поставьте под запятыми исходных чисел.
\end{enumerate}

Например,
\[
1{,}2+3{,}51
=
1{,}20+3{,}51
=
4{,}71.
\]

% =================================================

\sectiontitle{5. Вычитание десятичных дробей}

При вычитании используется тот же принцип:

\[
\important{\text{запятая под запятой}}.
\]

Например:
\[
2{,}410-1{,}159.
\]

\[
\begin{array}{r}
  2{,}410\\
- 1{,}159\\
\hline
  1{,}251
\end{array}
\]

Поэтому:
\[
\boxed{2{,}410-1{,}159=1{,}251}.
\]

Более сложный пример:
\[
10{,}12-3{,}453.
\]

Сначала допишем недостающий нуль:
\[
10{,}120-3{,}453.
\]

Тогда:
\[
\begin{array}{r}
 10{,}120\\
- 3{,}453\\
\hline
  6{,}667
\end{array}
\]

Следовательно,
\[
\boxed{10{,}12-3{,}453=6{,}667}.
\]

Если одно из чисел целое, его можно представить как десятичную дробь:

\[
10=10{,}00.
\]

Например:
\[
10-1{,}12
=
10{,}00-1{,}12
=
8{,}88.
\]

\important{Важно:}
при вычитании порядок чисел менять нельзя.
Уменьшаемое и вычитаемое определяются исходным выражением.

% =================================================

\sectiontitle{6. Умножение десятичных дробей}

При умножении правила отличаются от сложения и вычитания.

\important{Запятые при записи столбиком выравнивать не требуется.}

Удобно:

\begin{enumerate}
  \item временно не учитывать запятые;
  \item выполнить обычное умножение натуральных чисел;
  \item посчитать общее количество цифр после запятой
        в обоих множителях;
  \item отделить в произведении справа столько же цифр.
\end{enumerate}

Число с большим количеством цифр обычно удобно записывать сверху.

\subsection*{\color{accent}Пример 1}

\[
1{,}5\cdot3=4{,}5.
\]

У первого множителя одна цифра после запятой,
у второго --- ни одной.

Всего:
\[
1+0=1.
\]

Поэтому в произведении отделяется одна цифра справа.

\subsection*{\color{accent}Пример 2}

\[
1{,}5\cdot3{,}2.
\]

Временно вычисляем:
\[
15\cdot32=480.
\]

После запятой в исходных множителях всего две цифры:

\[
1+1=2.
\]

Получаем:
\[
4{,}80=4{,}8.
\]

Следовательно,
\[
\boxed{1{,}5\cdot3{,}2=4{,}8}.
\]

\subsection*{\color{accent}Пример 3}

\[
3{,}12\cdot1{,}5.
\]

Временно:
\[
312\cdot15=4680.
\]

Количество цифр после запятой:
\[
2+1=3.
\]

Значит:
\[
4{,}680=4{,}68.
\]

Итак,
\[
\boxed{3{,}12\cdot1{,}5=4{,}68}.
\]

\subsection*{\color{accent}Если цифр не хватает}

Например:
\[
0{,}3\cdot0{,}2.
\]

Сначала:
\[
3\cdot2=6.
\]

Нужно отделить две цифры справа.
Поэтому дописываем нуль перед шестёркой:

\[
\boxed{0{,}3\cdot0{,}2=0{,}06}.
\]

% =================================================

\sectiontitle{7. Деление на натуральное число}

При делении десятичной дроби на натуральное число удобно
использовать обычное деление столбиком.

\important{Делитель обязательно должен быть отличен от нуля.}

То есть выражение
\[
a:b
\]
имеет смысл только при
\[
b\ne0.
\]

\subsection*{\color{accent}Алгоритм}

\begin{enumerate}
  \item Начинайте деление как обычно.
  \item Когда в делимом переходите через запятую,
        поставьте запятую в частном.
  \item Продолжайте деление цифр дробной части.
  \item Если цифры закончились, а остаток остался,
        можно дописывать справа нули.
\end{enumerate}

Например:
\[
15{,}6:3=5{,}2.
\]

Проверка:
\[
5{,}2\cdot3=15{,}6.
\]

Ещё один пример:
\[
17{,}5:5=3{,}5.
\]

Проверка:
\[
3{,}5\cdot5=17{,}5.
\]

Также:
\[
2{,}4:4=0{,}6.
\]

% =================================================

\sectiontitle{8. Когда при делении целых чисел появляется десятичная дробь}

Рассмотрим:
\[
25:4.
\]

Имеем:
\[
4\cdot6=24.
\]

После первого шага остаётся:
\[
25-24=1.
\]

Целая часть закончилась, поэтому в частном ставим запятую
и дописываем нуль:

\[
1\longrightarrow10.
\]

Далее:
\[
10:4=2
\]
с остатком \(2\).

Снова дописываем нуль:
\[
20:4=5.
\]

Получаем:
\[
\boxed{25:4=6{,}25}.
\]

\important{Запятая в частном ставится один раз.}
После неё при необходимости можно продолжать дописывать нули
к делимому.

Проверка:
\[
6{,}25\cdot4=25.
\]

% =================================================

\sectiontitle{9. Среднее арифметическое}

\important{Среднее арифметическое нескольких чисел} ---
это сумма всех чисел, делённая на их количество.

Для чисел
\[
a_1,a_2,\ldots,a_n,
\qquad n\ge1,
\]
среднее арифметическое обозначают, например, через
\(\overline a\):

\[
\boxed{
\overline a=
\frac{a_1+a_2+\cdots+a_n}{n}
}.
\]

\subsection*{\color{accent}Пример}

Найдём среднее арифметическое чисел
\[
4,\ 5,\ 6,\ 10,\ 12.
\]

Сначала вычисляем сумму:
\[
4+5+6+10+12=37.
\]

Чисел всего пять, поэтому:
\[
\overline a=\frac{37}{5}=7{,}4.
\]

Ответ:
\[
\boxed{7{,}4}.
\]

\subsection*{\color{accent}Удобная группировка}

Для чисел
\[
4{,}9,\quad5{,}1,\quad5,\quad4{,}8,\quad5{,}2
\]
можно заметить:

\[
4{,}9+5{,}1=10,
\]

\[
4{,}8+5{,}2=10.
\]

Поэтому:
\[
10+10+5=25.
\]

Следовательно:
\[
\overline a=\frac{25}{5}=5.
\]

Такое удобное объединение слагаемых помогает выполнять
вычисления быстрее.

% =================================================

\sectiontitle{10. Простейшие уравнения}

При решении уравнений используются
\important{обратные арифметические действия}.

\subsection*{\color{accent}Основные формы}

\renewcommand{\arraystretch}{1.4}

\begin{center}
\begin{tabularx}{0.88\linewidth}{
  >{\raggedright\arraybackslash}X
  >{\raggedright\arraybackslash}X
}
\toprule
Уравнение & Неизвестное \\
\midrule
\(x+a=b\) & \(x=b-a\) \\

\(x-a=b\) & \(x=b+a\) \\

\(ax=b,\quad a\ne0\) & \(x=\dfrac{b}{a}\) \\

\(\dfrac{x}{a}=b,\quad a\ne0\) & \(x=ab\) \\
\bottomrule
\end{tabularx}
\end{center}

\subsection*{\color{accent}Примеры}

\[
2x=10
\]

\[
x=10:2=5.
\]

\vspace{0.4em}

\[
x+2=10
\]

\[
x=10-2=8.
\]

\vspace{0.4em}

\[
x-2=10
\]

\[
x=10+2=12.
\]

\subsection*{\color{accent}Уравнение со скобками}

\[
2(x-8)=10.
\]

Сначала делим обе части на \(2\):

\[
x-8=5.
\]

Затем прибавляем \(8\):

\[
\boxed{x=13}.
\]

Проверка:
\[
2(13-8)=2\cdot5=10.
\]

\subsection*{\color{accent}Уравнение с делением}

\[
\frac{x-8}{2}=10.
\]

Умножаем обе части на \(2\):

\[
x-8=20.
\]

Прибавляем \(8\):

\[
\boxed{x=28}.
\]

Проверка:
\[
\frac{28-8}{2}
=
\frac{20}{2}
=
10.
\]

\important{Правило:}
одно и то же допустимое действие выполняется с обеими частями
уравнения.

% =================================================

\sectiontitle{11. Что обязательно проверять}

\begin{itemize}
  \checkmarkitem
  При сложении и вычитании десятичных дробей
  запятые стоят строго друг под другом.

  \checkmarkitem
  Недостающие цифры дробной части при сложении и вычитании
  можно заменить нулями справа.

  \checkmarkitem
  Нули внутри десятичной записи удалять нельзя:
  \(4{,}08\ne4{,}8\).

  \checkmarkitem
  При умножении подсчитано \textbf{общее} количество цифр
  после запятых в обоих множителях.

  \checkmarkitem
  Если при постановке запятой в произведении цифр недостаточно,
  слева дописаны нули.

  \checkmarkitem
  При делении запятая не потеряна.

  \checkmarkitem
  Делитель не равен нулю.

  \checkmarkitem
  При нахождении среднего арифметического сумма разделена
  именно на \textbf{количество чисел}.

  \checkmarkitem
  При решении уравнения преобразование выполнено
  с обеими его частями.

  \checkmarkitem
  Полученный ответ проверен подстановкой или обратным действием.
\end{itemize}

% =================================================
% ТРЕНИРОВОЧНЫЙ БЛОК
% =================================================

\clearpage

\section*{\color{darkaccent}Тренировочный блок · Путь В}

Задания концептуально повторяют разобранные приёмы.
Решите их самостоятельно. Решения в этом разделе не приводятся.

\vspace{0.7em}

\subsection*{\color{accent}Упражнение 1. Вычисления и порядок действий}

Вычислите:

\begin{enumerate}[label=\arabic*)]
  \item \(9000-4687\);
  \item \(840-7\cdot(18+22)\).
\end{enumerate}

\vspace{0.7em}

\subsection*{\color{accent}Упражнение 2. Сравнение десятичных дробей}

Поставьте вместо \(\square\) знак \(<\), \(>\) или \(=\):

\begin{enumerate}[label=\arabic*)]
  \item \(5{,}4\ \square\ 5{,}400\);
  \item \(4{,}08\ \square\ 4{,}008\);
  \item \(1{,}245\ \square\ 1{,}254\);
  \item \(12{,}03\ \square\ 9{,}999\).
\end{enumerate}

\vspace{0.7em}

\subsection*{\color{accent}Упражнение 3. Действия с десятичными дробями}

Вычислите:

\begin{enumerate}[label=\arabic*)]
  \item \(3{,}47+1{,}286\);
  \item \(10{,}04-2{,}735\);
  \item \(2{,}4\cdot1{,}35\);
  \item \(18{,}9:7\).
\end{enumerate}

\vspace{0.7em}

\subsection*{\color{accent}Упражнение 4. Среднее арифметическое}

Найдите среднее арифметическое чисел

\[
4{,}7,\qquad
5{,}3,\qquad
6,\qquad
4{,}9,\qquad
5{,}1.
\]

\vspace{0.7em}

\subsection*{\color{accent}Упражнение 5. Уравнения}

Решите уравнения:

\begin{enumerate}[label=\arabic*)]
  \item \(3(x-4)=15\);
  \item \(\dfrac{x-6}{3}=7\).
\end{enumerate}

% =================================================
% ОТВЕТЫ
% =================================================

\clearpage

\section*{\color{darkaccent}Ответы к тренировочному блоку}

\renewcommand{\arraystretch}{1.45}

\begin{table}[ht]
\centering
\caption{Ответы}
\begin{tabularx}{\linewidth}{
  >{\bfseries}p{0.14\linewidth}
  >{\raggedright\arraybackslash}X
}
\toprule
Упражнение & Ответ \\
\midrule

1 &
1) \(4313\);
\quad
2) \(560\).
\\

2 &
1) \(=\);
\quad
2) \(>\);
\quad
3) \(<\);
\quad
4) \(>\).
\\

3 &
1) \(4{,}756\);
\quad
2) \(7{,}305\);
\quad
3) \(3{,}24\);
\quad
4) \(2{,}7\).
\\

4 &
\(5{,}2\).
\\

5 &
1) \(x=9\);
\quad
2) \(x=27\).
\\

\bottomrule
\end{tabularx}
\end{table}

\vfill

\begin{center}
\begin{tikzpicture}[x=1cm,y=1cm]
  \draw[
    accent!45,
    line width=0.9pt
  ]
  (0,0)
  .. controls (2.4,0.4) and (4.6,0.4) .. (7,0)
  .. controls (9.4,-0.4) and (11.6,-0.4) .. (14,0);
\end{tikzpicture}
\end{center}

\end{document}